\BeleskeID{04-s014}
% element: 04-s014:e01 — Vp=Vq i mogućnost prelaska preko decimalnog sistema bez obaveze
% element: 04-s014:e02 — Primer 236 u osnovi 7 u binarni sistem
% element: 04-s014:e03 — Tri koraka deljenja, delovi broja, ostaci i količnik 116
% element: 04-s014:e04 — Navika na decimalne tablice i potreba ustanovljavanja algoritama
% element: 04-s014:d01

Decimalne tablice sabiranja i množenja toliko su uvežbane da često ne primećujemo algoritam. Da bismo računanje opisali digitalnom sistemu, potrebno je izdvojiti korake i pravilo prenosa ostatka; isti princip može da se primeni i u drugoj osnovi.

Pri pisanom deljenju $236_7$ sa $2$, posle deljenja prve cifre ostaje $0$. Spuštanjem cifre $3$ obrađujemo $03_7$, pa dobijamo količnik $1$ i ostatak $1$. Taj ostatak se prenosi u sledeće poziciono mesto: spuštanjem cifre $6$ nastaje $16_7$, čija je vrednost $1\cdot7+6=13_{10}$. Zato je sledeća cifra količnika $6$, sa ostatkom $1$.
\[236_7=2\cdot116_7+1.\]
Broj $16$ u tom koraku nije decimalni šesnaest. Deljenje se izvodi u osnovi $7$, pa i delovi broja i količnik koriste tu osnovu. Decimalni račun može služiti proveri, ali nije obavezna međufaza konverzije.
